\section{先从磁盘读回 Stage 1}

\topic{磁盘布局先于读取代码确定}
项目不采用 MBR 或文件系统作为最早入口。LBA 0 保存一扇区自描述头，负载从
头中声明的 LBA 开始，当前构建值为 1。宿主镜像工具与 ROM 固件分别实现同一
字段契约，使构建产物和目标解析能够相互校验。

\begin{osgridlongtable}{L{0.12\textwidth}L{0.13\textwidth}L{0.25\textwidth}L{0.38\textwidth}}
\textbf{偏移} & \textbf{宽度} & \textbf{字段} & \textbf{固件约束}\\ \hline
\endfirsthead
\textbf{偏移} & \textbf{宽度} & \textbf{字段} & \textbf{固件约束}\\ \hline
\endhead
\code{0x00} & 8 字节 & magic & 精确等于 \code{OSSTAGE1}\\ \hline
\code{0x08} & 2 字节 & version & 当前等于 1\\ \hline
\code{0x0A} & 2 字节 & header size & 当前等于 28\\ \hline
\code{0x0C} & 2 字节 & load segment & 形成合法 RAM 加载窗口\\ \hline
\code{0x0E} & 2 字节 & entry offset & 小于补零后负载长度\\ \hline
\code{0x10} & 2 字节 & sector count & 1 到 64\\ \hline
\code{0x12} & 2 字节 & flags & 当前必须为零\\ \hline
\code{0x14} & 4 字节 & payload LBA & LBA28 且不越过镜像\\ \hline
\code{0x18} & 2 字节 & payload checksum & 等于负载 16 位字和\\ \hline
\code{0x1A} & 2 字节 & header checksum & 使描述符整扇区字和为零\\ \hline
\end{osgridlongtable}

描述符剩余字节同样参加校验。这样保留区被意外写入也会失败，未来只有提升格式
版本并定义新字段后才能使用。magic 识别格式，校验发现变化，范围检查保护内存，
三者承担不同职责。

\topic{加载窗口避开固件临时状态}
描述符读入物理地址 \code{0x0500..0x06FF}，实模式栈从 \code{0x7000}
向低地址增长。Stage 1 当前装入 \code{0x0800:0x0000}，即物理
\code{0x8000}。固件只接受：
\[
  0x8000 \leq \mathrm{load_begin}
  < \mathrm{load_end} \leq 0x9FC00.
\]
结束地址由扇区数乘 512 后与起始地址相加，乘法、加法和 LBA 结束位置都先在
32 位寄存器中检查。最多 64 个扇区同时保证 \code{rep insw} 推进 16 位 DI
时不会绕回。

\topic{ATA PIO 状态机逐扇区提交}
固件选择主 IDE 通道主盘，把 LBA 的低、中、高字节写入
\code{0x1F3..0x1F5}，高四位与设备选择写入 \code{0x1F6}，再向
\code{0x1F7} 写 \code{0x20}。四次读取 alternate status 提供设备选择后的
400 ns 延时。

命令发出后最多读取状态 \code{0xFFFF} 次。BSY 仍为一时继续等待；ERR 或 DF
出现时立即返回设备错误；只有 BSY 清零且 DRQ 置位才从 \code{0x1F0} 执行
256 次 16 位读取。每条命令只读一个扇区，完成后由外层增加 LBA 并减少剩余
计数，使错误位置保持精确。

\topic{校验完成后才交出指令指针}
所有扇区进入 RAM 后，固件把补零负载解释为小端 16 位字序列，计算模
\(2^{16}\) 的和并与描述符比较。成功日志在控制转移前输出
\code{STAGE1\_LOADED}。随后依次压入目标 CS 和 IP，执行远返回；处理器重载
CS 隐藏状态，开始在 \code{0x0800:0x0000} 取指。

Stage 1 不能假定 DS 随 CS 改变。其第一段代码用 \code{push cs}/
\code{pop ds} 建立自身常量地址，再通过已经初始化的 COM1 独立输出
\code{[OS][STAGE1] ENTERED}。这条标记证明磁盘读取、校验、RAM 写入、远跳转
和新代码自身寻址全部完成。

\topic{失败证据对应不同边界}
\small
\begin{osgridlongtable}{L{0.36\textwidth}L{0.23\textwidth}L{0.27\textwidth}}
\textbf{标记} & \textbf{失败边界} & \textbf{测试输入}\\ \hline
\code{IDE\_TIMEOUT} & 状态始终忙或始终无 DRQ & 同源 ROM 强制 BSY\\ \hline
\code{IDE\_ERROR} & 设备返回 ERR/DF & 同源 ROM 强制 ERR\\ \hline
\code{STAGE1\_HEADER\_INVALID} & 格式、范围或描述符校验 & 损坏 magic\\ \hline
\code{STAGE1\_CHECKSUM\_INVALID} & 已读取负载内容变化 & 翻转负载字节\\ \hline
\end{osgridlongtable}
\normalsize

每条失败路径禁止输出 \code{STAGE1\_LOADED} 与 Stage 1 进入标记。宿主单元
测试另行拒绝短镜像和越界 LBA；固定种子测试覆盖 256 组有效负载往返和 256
组越界输入，使格式性质不依赖单一样例。

\section{把 Kernel ELF 放到不会互相覆盖的位置}

\topic{磁盘位置不能成为构建系统的秘密}
生成 \filepath{kernel.elf} 以后，Stage 1 仍需要知道从哪个 LBA 读取多少
字节。如果构建脚本把内核简单拼接在 Stage 1 之后，Stage 1 每增长一个扇区，
Kernel 的位置就随之变化；目标代码要么硬编码一个会漂移的数字，要么依赖宿主
额外注入的隐含常量。二者都破坏了磁盘自身可以解释的原则。

项目因而保留 v0.2 的 Stage 1 格式，并在同一原始磁盘上增加第二套独立容器：

\begin{osgridlongtable}{L{0.16\textwidth}L{0.29\textwidth}L{0.43\textwidth}}
\textbf{LBA} & \textbf{内容} & \textbf{边界}\\ \hline
\endfirsthead
\textbf{LBA} & \textbf{内容} & \textbf{边界}\\ \hline
\endhead
0 & Stage 1 描述符 & 固件最先读取并验证\\ \hline
1 到 64 & Stage 1 最大负载区 & v1 格式最多允许 64 个扇区\\ \hline
65 & Kernel v1 描述符 & 固定位置，不随 Stage 1 当前长度漂移\\ \hline
66 起 & 精确 \filepath{kernel.elf} & 文件后只补零到扇区边界\\ \hline
\end{osgridlongtable}

固定 LBA 65 并不浪费重要容量：v0.4 的 2 MiB 教学磁盘已包含 4096 个扇区，
自 v1.6 起逻辑 1 GiB 镜像空间更充足；无论容量如何，稳定
交接比节省几十个空扇区更重要。宿主打包器先完整解析 Stage 1 描述符，再证明
其结束 LBA 不大于 65；它不因“当前 Stage 1 只有七个扇区”就放弃最大边界
检查。

\topic{Kernel 描述符使用可扩展宽度但服从当前硬件}\par
\begin{osgridlongtable}{L{0.13\textwidth}L{0.12\textwidth}L{0.24\textwidth}L{0.39\textwidth}}
\textbf{偏移} & \textbf{宽度} & \textbf{字段} & \textbf{当前约束}\\ \hline
\endfirsthead
\textbf{偏移} & \textbf{宽度} & \textbf{字段} & \textbf{当前约束}\\ \hline
\endhead
\code{0x00} & 8 & magic & 精确等于 \code{OSKERN64}\\ \hline
\code{0x08} & 2 & version & 等于 1\\ \hline
\code{0x0A} & 2 & header size & 等于 48\\ \hline
\code{0x0C} & 4 & flags & 当前必须为零\\ \hline
\code{0x10} & 8 & payload LBA & 当前等于 66，且不得超过 LBA28\\ \hline
\code{0x18} & 8 & file size & 精确 ELF 文件字节数\\ \hline
\code{0x20} & 8 & sector count & 文件长度向上取整到 512 字节\\ \hline
\code{0x28} & 4 & payload CRC32 & 只覆盖精确文件，不覆盖尾部补零\\ \hline
\code{0x2C} & 4 & descriptor CRC32 & 计算时本字段置零，覆盖整扇区\\ \hline
\code{0x30..0x1FF} & 464 & reserved & 当前版本必须全零\\ \hline
\end{osgridlongtable}

LBA、文件长度和扇区数使用 64 位字段，是为了避免格式很快因容量增长而升级；
这不表示当前 ATA 驱动已经支持 LBA48。描述符解析仍把 LBA 与
\code{0x0FFFFFFF} 比较。格式的表达范围、驱动的协议范围和真实磁盘范围是三
个不同约束，必须分别成立。

保留区全零规则也不能被 CRC 取代。CRC 能发现位变化，但攻击者或错误工具可以
在修改保留字节后重新计算 CRC；版本 1 解析器若只检查校验值，就会默许未定义
语义。先检查“当前字段是否允许”，再检查“字节是否完整”，两者回答的问题
不同。

\topic{CRC32 检测损坏但不宣称提供认证}
Stage 1 v1 使用 16 位字和，计算极小且适合最早 ROM。ELF 文件更大，v0.4
选择 CRC32 检测常见的突发位错误、错位和内容替换。描述符 CRC 覆盖完整
512 字节，Kernel CRC 只覆盖 \code{file size} 指定的精确字节；否则扇区尾部
补零也会意外成为文件内容的一部分。

CRC32 不是密码学散列，也没有密钥。能够主动同时改写文件和描述符的人可以
重新计算 CRC。因此这里的安全结论是“偶然损坏可检测”，不是“来源可信”。
ELF 的长度、对齐、权限、重叠和入口检查仍然不可省略。

\topic{验证顺序保持先证明再写入}
宿主组合磁盘时按以下顺序提交：

\begin{enumerate}[label=\arabic*.]
  \item 完整验证 Stage 1 镜像和它的负载结束位置。
  \item 完整解析 Kernel ELF64，确认入口和 PT\_LOAD 契约。
  \item 计算精确文件长度、扇区数与两个 CRC32。
  \item 证明 LBA 65 之后的目标区域仍全零且落在磁盘内。
  \item 最后一次性写入描述符与 ELF 字节。
\end{enumerate}

解析组合磁盘时按相反顺序检查这些对象：Stage 1、阶段分离、Kernel 描述符、
磁盘范围、文件 CRC、ELF64。单元测试覆盖重叠、保留区和损坏输入；固定种子
测试完成 128 组不同文件长度往返和 256 组随机文件字节破坏；集成测试审计
真实 \filepath{boot\_disk.img}，同时要求两份故障镜像被拒绝。

这一层没有让 Python 替代目标机工作。宿主工具只是生成和独立审计输入；
64 位 Stage 1 已自行通过 ATA PIO 读取 LBA 65，执行 CRC32 与 ELF 程序头
验证，再把段复制到 \code{0x00100000}。两端实现同一格式，才能互相发现错误。

\topic{为什么映射从 2 MiB 扩展到 64 MiB}
v0.3 只需要让正在执行的 Stage 1、页表和早期栈跨过分页开启点，首个 2 MiB
大页已经足够。v0.4 增加了三个新对象：低端 ELF 暂存文件、1 MiB 起的加载段
以及接近 64 MiB 的独立内核栈。若仍只映射第一张大页，代码段或栈的第一次
访问就会触发页故障；IDT 尚未建立时，页故障会升级为双重故障甚至三重故障，
表现为机器突然复位。

项目没有因此增加四级页表数量。一个 PD 含 512 个 64 位项，每个设置 PS 的
项映射 2 MiB；顺序填写前 32 项即可覆盖：
\[
  32\times 2\ \mathrm{MiB}=64\ \mathrm{MiB}.
\]
每项的物理基址是索引乘 2 MiB，再与 present、writable 和 page-size 位组合。
构造后逐项回读，而不是只检查第一项。这既验证循环边界，也防止索引或步长错误
让高地址栈成为未映射洞。

\topic{物理区间表同时是所有权表}
\small
\begin{osgridlongtable}{L{0.27\textwidth}L{0.27\textwidth}L{0.34\textwidth}}
\textbf{半开区间} & \textbf{所有者} & \textbf{为何放在这里}\\ \hline
\endfirsthead
\textbf{半开区间} & \textbf{所有者} & \textbf{为何放在这里}\\ \hline
\endhead
\code{0x00008000..0x00010000} & Stage 1 & 与实模式加载位置一致，最大 32 KiB\\ \hline
\code{0x00010000..0x00013000} & 页表 & 三张 4 KiB 对齐表\\ \hline
\code{0x00013000..0x00013200} & Kernel 描述符 & 一个完整磁盘扇区\\ \hline
\code{0x00014000..0x00014068} & BootInfo v2 & 独立页面内的 104 字节 ABI\\ \hline
\code{0x00015000..0x00016000} & 加载器临时状态 & 程序头遍历索引和计数\\ \hline
\code{0x00016000..0x00017000} & 内存图元数据 & 交接条目数量\\ \hline
\code{0x00017000..0x00018000} & \code{fw\_cfg} 暂存 & 当前 56 字节目录名称\\ \hline
\code{0x00018000..0x00018C00} & 物理内存图 & 最多 128 个 24 字节项\\ \hline
\code{0x00100000..0x03600000} & \code{PT\_LOAD} 窗口 & 与高地址源暂存区严格分离\\ \hline
\code{0x03600000..0x03E00000} & Kernel 文件暂存 & 经 E820 验证的 8 MiB 上限\\ \hline
\code{0x03FEF000..0x03FFF000} & 初始内核栈 & 64 KiB，向低地址增长\\ \hline
\end{osgridlongtable}
\normalsize

表中的空洞不是遗漏。描述符之后按页分开 BootInfo 和临时状态，使 GDB 可以
按稳定地址检查，也给结构增长留下明确边界。v0.10 曾把 Kernel 文件暂存放在
62--63 MiB；v1.9 因载荷超过 1 MiB，改为 54--62 MiB 的 8 MiB 区间。
Stage 1 在读盘前要求完整区间落入同一个 E820 usable 条目，不能只凭
\code{-m 64} 猜测可写。加载目标结束于 \code{0x03600000}，因此两遍加载
第二遍复制时不可能覆盖仍在读取的源 ELF。

区间全部采用半开语义。验证段结束地址时先证明
\(\mathrm{begin}+\mathrm{size}\) 不发生 64 位溢出，再比较结束地址不超过
窗口上界；不能先做可能回绕的加法，然后用回绕后的“小数”通过检查。

\topic{寄存器变宽不会让端口字段变宽}
CPU 进入 64 位模式后，通用寄存器宽度变成 64 位，但 ATA LBA28 命令块仍由
8 位地址寄存器和 16 位数据端口组成。Stage 1 把 32 位 LBA 拆成低、中、高
字节和高四位，选择 primary master，发送 \code{READ SECTORS}，每次只读
一个扇区。数据目的地址放在 RDI，\code{rep insw} 仍执行 256 次 16 位输入。

每扇区单独提交看似不如批量高效，却让三个边界非常清晰：扇区计数字段永远为
1，不遇到“零表示 256”的历史语义；失败时当前 LBA 明确；目标指针每次只前进
512 字节，容易证明不会越过暂存区。性能优化应在正确状态机之上进行，而不是
让最早加载器先承担批量命令的额外状态。

\topic{一次真实的超时误判}
Stage 1 尚无中断时钟和任务调度器，只能用固定次数预算防止设备永久忙。状态
判定顺序保持：
\[
\begin{array}{ll}
\mathrm{BSY}=1 & \rightarrow\ \text{继续等待}\\
\mathrm{BSY}=0,\ \mathrm{ERR}\lor\mathrm{DF}=1
  & \rightarrow\ \text{设备错误}\\
\mathrm{BSY}=0,\ \mathrm{DRQ}=0 & \rightarrow\ \text{继续等待}\\
\mathrm{BSY}=0,\ \mathrm{DRQ}=1 & \rightarrow\ \text{读取数据}.
\end{array}
\]

Git 历史保留了一次具体修改。Kernel 增长以后，Stage 1 要提交上千个单扇区
事务；宿主繁忙、TCG 线程被抢占时，原来的 \code{0x0000FFFF} 次预算偶尔会
把正常设备判成超时。串口最后出现的是
\code{[OS][STAGE1] KERNEL_ATA_TIMEOUT}，而不是设备 ERR，调查因而落在轮询
预算和宿主调度之间。

提交 \code{c588157} 把
\filepath{source/boot/stage1/include/kernel\_loader.inc} 中的预算提高到
\code{0x000FFFFF}。失败镜像仍然让状态永久保持 BSY，确认循环最终会退出。
这个改动没有把固定次数变成稳定墙钟；进入拥有单调时钟的内核以后，等待改用
deadline，Stage 1 外层的 QEMU 测试也继续保留独立墙钟上限。

超时与设备错误分别映射为 \code{KERNEL\_ATA\_TIMEOUT} 和
\code{KERNEL\_ATA\_ERROR}。测试变体在真实轮询函数内强制状态，而不是在
上层直接打印失败文字，因此能够证明错误分支、日志收敛和后续里程碑禁止条件。

\topic{为什么它比简单字和更适合 Kernel}
Stage 1 的 16 位字和能发现许多单字变化，但交换两个字不会改变和，某些成对
误差也会抵消。CRC 把比特串解释为 \(GF(2)\) 上的多项式，除法中的加减都由
异或实现；余数对连续突发错误有更系统的检测能力。IEEE CRC32 的标准多项式
常写为 \code{0x04C11DB7}。按最低位先处理时使用它的反射形式
\code{0xEDB88320}。

项目算法以 \code{0xFFFFFFFF} 初始化余数，每个输入字节先与余数低字节异或，
再进行八轮：
\[
  c_{i+1}=
  \begin{cases}
    (c_i\gg1)\oplus\mathtt{EDB88320},&c_i\text{ 的最低位为 1},\\
    c_i\gg1,&\text{否则}.
  \end{cases}
\]
所有字节处理后逐位取反。逐位实现比 256 项查表慢，但代码和常量更少，算法
路径可以直接与宿主标准库结果对照；当前最大 8 MiB 文件仍有明确上限。

\topic{描述符、自身字段和尾部填充}
描述符 CRC 覆盖整个 512 字节，但 CRC 字段不可能同时包含自己最终值并直接
参与普通计算。目标实现先保存磁盘中的字段，把内存副本对应四字节置零，计算
后恢复，再比较结果。恢复动作很重要：后续 BootInfo 和调试器看到的描述符仍
应等于磁盘内容。

负载 CRC 只覆盖 \code{file size} 个精确字节。扇区中最后的补零不属于 ELF，
但也不能完全忽略：Stage 1 另行扫描从精确文件末尾到
\(\mathrm{sector_count}\times512\) 的区域，要求全部为零。这样既保持 CRC
语义稳定，也不允许未校验数据藏在容器尾部。

\section{先检查 ELF，再写入目标内存}

\topic{第一层从文件身份建立边界}
ELF 头至少 64 字节。加载器先检查 magic、64 位类别、小端序、当前标识版本、
\code{ET\_EXEC}、x86-64 机器号、ELF 版本、固定入口、头大小和程序头项
大小。程序头数量必须位于 1 到 64，程序头表偏移不得位于头之前，表末尾必须
落在精确文件范围。

最多 64 项是当前 Stage 1 的资源预算；ELF 标准本身允许更多。这个上限为两两重叠
检查给出最多 \(64\times63/2=2016\) 次区间比较，也避免损坏文件用极大计数
拖长最早启动。格式字段较宽不意味着实现必须无条件接受所有可表达值。

\topic{每个加载段有局部不变量}
程序头类型不是 \code{PT\_LOAD} 时可以跳过；加载段则必须全部满足：

\begin{enumerate}[label=\arabic*.]
  \item 权限只含 R、W、X，且 R 必须存在；
  \item W 与 X 不能同时出现；
  \item \code{p\_filesz <= p\_memsz}，且内存长度非零；
  \item 文件偏移加文件长度不溢出且不越过精确 ELF；
  \item 对齐字段精确为 4 KiB，文件偏移与虚拟地址页内偏移相同；
  \item 虚拟地址等于物理地址，暂不引入重定位或非恒等目标；
  \item 内存区间完整落入 \code{0x00100000..0x03600000}。
\end{enumerate}

要求 R 权限和 W\^X 不是 Stage 1 大页已经能强制出来的最终安全策略；初始
2 MiB 大页仍以较粗权限映射。它先保证 ELF 自身给出的意图没有矛盾；v0.6
内核随后依据链接器段边界拆成 4 KiB 页并真正执行权限。

\topic{跨段不变量不能在单个程序头中检查}
单个段合法不代表段集合合法。加载器把当前段与所有先前
\code{PT\_LOAD} 做半开区间比较，任何重叠都拒绝。否则后复制的段会覆盖先
复制的段，最终内存内容依赖程序头顺序。固定入口
\code{0x00100000} 还必须位于一个带 X 的加载段，不能仅等于 ELF 头中的
\code{e\_entry} 就通过。

第一遍只读取暂存文件和一页临时状态，不写
\code{0x00100000} 之后的目标区。等全部程序头通过，第二遍才对每个加载段：
\[
  \operatorname{copy}(p\_offset,p\_vaddr,p\_filesz),\qquad
  \operatorname{zero}(p\_vaddr+p\_filesz,p\_memsz-p\_filesz).
\]
验证失败因此保持“没有半个内核”的状态，日志和 GDB 现场仍可解释。

\topic{BSS 把语言保证变成加载器责任}
BSS 名称来自早期汇编工具中的 Block Started by Symbol。未显式初始化或零
初始化的全局对象不必在文件中存储一串零；链接器让段的内存大小大于文件大小，
把差值交给加载器清零。这能缩小磁盘文件，却把责任从编译器转移到了装载阶段。

C++ 规定静态存储期对象在进一步初始化前先零初始化。freestanding 环境不会
有宿主运行时替加载器弥补错误。本项目在 BSS 放置一个 64 位探针：内核入口
第一次读取必须为零，随后写入一个命名的非零值并输出
\code{BSS\_ZEROED}。若 Stage 1 只复制 \code{p\_filesz} 而忘记清零，
QEMU 系统测试将在第一次 C++ 交接中暴露问题。

\topic{没有外部 BIOS，就不能调用一个不存在的 E820 服务}
经典 PC 教材常让实模式代码执行 \code{INT 15h, EAX=E820h}。这个接口不是
CPU 指令自动提供的能力，而是 BIOS 安装的软件中断服务。本项目的 ROM 从复位
入口开始自行实现，没有 SeaBIOS；若为了获取内存图临时换回外部 BIOS，就等于
把完整启动链最关键的硬件发现步骤交给了别人的固件。

QEMU PC 另有 \code{fw\_cfg} 配置设备。PIO 形式只有两个端口：向
\code{0x510} 写 16 位 selector，再从 \code{0x511} 顺序读取字节。QEMU
模拟端口和数据，本项目仍负责协议状态机，边界与 UART、IDE 完全相同。

Stage 1 首先选择 key \code{0x0000}，读四字节签名并要求为 \code{QEMU}；
然后选择 key \code{0x0019} 读取文件目录。目录数量和每项的 size/selector
采用大端序，而 x86 端口读到的单字节不自动完成转换。汇编逐字节左移和合并，
把外部窄字段提升到 64 位后才比较上限。

\topic{目录流式遍历避免新的早期内存依赖}
每个目录项固定 64 字节：

\begin{osgridlongtable}{L{0.18\textwidth}L{0.19\textwidth}L{0.49\textwidth}}
\textbf{偏移} & \textbf{宽度} & \textbf{语义}\\ \hline
\endfirsthead
\textbf{偏移} & \textbf{宽度} & \textbf{语义}\\ \hline
\endhead
\code{0x00} & 32 bit & 文件 size，大端\\ \hline
\code{0x04} & 16 bit & 重新选择文件所需 selector，大端\\ \hline
\code{0x06} & 16 bit & reserved\\ \hline
\code{0x08} & 56 byte & NUL 结尾的文件名\\ \hline
\end{osgridlongtable}

Stage 1 不把整个目录复制到 RAM。它最多遍历 256 项，只把当前 56 字节名称
放到 \code{0x17000}，精确匹配 \code{etc/e820} 后立即停止。这样目录大小
不会迫使我们再设计一块可变暂存区，最坏循环次数也保持有界。

\topic{20 字节平台项怎样成为 24 字节内核 ABI}
\code{etc/e820} 的一个原始项是小端
\(\{u64\ base,u64\ length,u32\ type\}\)，共 20 字节。Stage 1 要求文件
长度非零、能被 20 整除且最多 128 项。读每项后在尾部补一个 32 位
attributes 零，形成 24 字节结构。扩展字段预留稳定
ABI 槽；内核当前不依据其作决策。

QEMU 当前给出的高端保留项可能排在低端 RAM 前面。文件顺序不是内核排序契约，
所以 Stage 1 在交接前按 base 做有界冒泡排序。内核随后仍检查单调和重叠；
加载器排序负责建立候选规范形状，内核验证负责保护自己的信任边界，两者不是
重复浪费。

签名、目录数量、名称、selector、长度或条目数任一不合法，只输出
\code{MEMORY\_MAP\_INVALID} 并停止。独立 QEMU 镜像在真实长模式边界强制这
条失败，验证它不会继续读取 Kernel 或构造 BootInfo。

\topic{十三个 64 位字段为何全部显式}
\small
\begin{osgridlongtable}{L{0.14\textwidth}L{0.33\textwidth}L{0.39\textwidth}}
\textbf{偏移} & \textbf{字段} & \textbf{当前值或来源}\\ \hline
\endfirsthead
\textbf{偏移} & \textbf{字段} & \textbf{当前值或来源}\\ \hline
\endhead
\code{0x00} & magic & 内存字节为 \code{OSBOOT64}\\ \hline
\code{0x08} & version & 2\\ \hline
\code{0x10} & structure size & 104 字节\\ \hline
\code{0x18} & Kernel 文件地址 & \code{0x03600000}\\ \hline
\code{0x20} & Kernel 文件长度 & 已通过 CRC 的精确长度\\ \hline
\code{0x28} & Kernel 入口 & \code{0x100000}\\ \hline
\code{0x30} & 加载段数量 & 第一遍验证所得，1 到 64\\ \hline
\code{0x38} & 页表根 & \code{0x10000}\\ \hline
\code{0x40} & 恒等映射大小 & \code{0x4000000}，即 64 MiB\\ \hline
\code{0x48} & 初始栈顶 & \code{0x3FFF000}\\ \hline
\code{0x50} & 物理内存图地址 & \code{0x18000}\\ \hline
\code{0x58} & 物理内存图条目数 & 已验证的 1 到 128\\ \hline
\code{0x60} & 内存图条目宽度 & 24 字节\\ \hline
\end{osgridlongtable}
\normalsize

这些字段统一使用 64 位，使跨阶段 ABI 不需要混合宽度，也给
后续扩展留出范围。磁盘字段、汇编偏移和 C++ 的 \code{uint64\_t} 必须三方
一致；不能用宿主平台的 \code{unsigned long} 或 \code{size\_t} 猜测宽度。
同时，大类型也不取消语义上限：文件仍最多 1 MiB，段数仍最多 64。

\topic{System V AMD64 ABI 规定第一次函数调用}
Stage 1 把 BootInfo 地址放入 RDI，把 RSP 设到独立栈顶，然后用
\code{CALL} 进入 \code{OsKernelEntry}。相较裸 \code{JMP}，\code{CALL}
自动压入返回地址，使入口栈状态符合普通 C++ 函数调用约定。入口被声明为
C ABI 和不返回，避免 C++ 名字修饰并明确生命周期；如果实现错误而返回，
Stage 1 仍有 \code{KERNEL\_RETURNED} 诊断边界。

初始阶段关闭红区，因为未来中断处理可能在当前栈指针下方写入状态，不能让
编译器假定那 128 字节永远安全。异常、RTTI、栈保护、线程安全局部静态对象
和宿主标准库也保持关闭，直到内核显式提供对应运行时。

\topic{重新初始化串口不是重复劳动}
Stage 1 已经配置 COM1，内核理论上可以直接发送。但那会形成一个没有写进 ABI
的隐藏依赖：任何固件改动都可能让内核日志失效。内核的
\code{SerialPort} 类用自己的固定宽度端口常量依次关闭中断、设置 DLAB、
写除数、恢复 8N1、配置 FIFO 与 modem control；每个字节发送仍有
\code{0xFFFF} 次 THRE 轮询预算。

类成员访问显式使用 \code{this->}，头文件只声明接口、源文件保存实现；常量
以项目、模块和功能为前缀。即使这是最早内核代码，也不降低工程命名和分层
要求。

\topic{每条成功日志排除了哪些故障}\par
\begin{osgridlongtable}{L{0.38\textwidth}L{0.48\textwidth}}
\textbf{标记} & \textbf{它实际证明的事实}\\ \hline
\endfirsthead
\textbf{标记} & \textbf{它实际证明的事实}\\ \hline
\endhead
\code{KERNEL\_HEADER\_VALID} & 目标 ATA 读取和描述符字段/CRC 成功\\ \hline
\code{KERNEL\_PAYLOAD\_VALID} & 精确 ELF CRC 和扇区补零成功\\ \hline
\code{KERNEL\_ELF\_VALID} & 第一遍 ELF 与全部段关系成功\\ \hline
\code{KERNEL\_SEGMENTS\_LOADED} & 第二遍复制与 BSS 清零完成\\ \hline
\code{BOOT\_INFO\_READY} & 十三字段 BootInfo v2 已写完\\ \hline
\code{KERNEL\_TRANSFER} & RDI、RSP 已准备，即将调用入口\\ \hline
\code{ENTERED} & C++ 入口和独立 COM1 代码确实执行\\ \hline
\code{BOOT\_INFO\_VALID} & 内核逐字段接受交接 ABI\\ \hline
\code{BSS\_ZEROED} & 加载器满足 C++ 零初始化前提\\ \hline
\code{CR3\_VALID} & 处理器实际页表根与交接事实一致\\ \hline
\code{MEMORY\_MAP\_VALID} & 内核接受排序、无重叠且无溢出的物理内存图\\ \hline
\code{PAGING\_READY} & 内核四级 4 KiB 页表已经激活\\ \hline
\code{HEAP\_SELF\_TEST\_PASSED} & 高半区堆真实写入、释放、合并与一致性检查成功\\ \hline
\code{USER\_RETURNED\_TO\_KERNEL} & 用户 exit 或异常已恢复原内核调用链\\ \hline
\code{READY} & 当前 v1.0 的入口、异常、内存、设备、进程、IPC、文件系统与交互环境自检全部完成\\ \hline
\end{osgridlongtable}

每扇区、每字节和每次轮询都不打印，避免串口耗时反过来改变启动行为。QEMU
宿主读取线程给每行加 \code{[QEMU][T+...ms]} 单调时间，这只表示宿主观察到
日志的时刻；来宾此时尚未建立 IRQ0 软件滴答，所以日志中没有来宾时间。

\topic{四层测试分别会抓到什么错误}
单元测试快速验证 BootInfo 每个错误状态、ELF 字段边界、CRC 与尾部补零；
集成测试检查真实 ELF、磁盘和内存区间；固定种子随机测试覆盖大量标识、目标
地址、文件长度、负载内容和填充变化；QEMU 系统测试才负责证明从
\code{0xFFFFFFF0} 到 C++ \code{READY} 的真实指令链。

失败镜像不是简单破坏后让宿主拒绝。语义非法 ELF 会重新计算描述符与负载
CRC，使目标先通过两层校验，再在 ELF 语义处失败；ATA 超时和 ERR
变体则直接作用于 Stage 1 的轮询函数。成功路径与五类 Kernel 失败路径共同
保证日志没有把错误归入笼统的“启动失败”。

完整测试命令会运行构建图自动登记的全部 CTest。代码量统计不包含测试程序、
宿主 Python 工具和教材内容，只覆盖 \filepath{source/} 中实际进入目标系统的
C++、头文件和 Intel 语法汇编。这样统计结果描述的是机器上真实执行的实现，
不会随测试数量或文字篇幅变化。
